\chapter{Kritische Bewertung und Lösungsansatz}

\section{Kritische Bewertung der bestehenden Ansätze}
\label{ch:kritische-bewertung}

Die im vorigen Kapitel vorgestellten regulatorischen und 
technischen Ansätze stellen insgesamt einen erheblichen 
Fortschritt dar, sie lösen das Souveränitätsproblem jedoch 
nicht vollständig.

Der Digital Markets Act, der Data Act und verwandte 
Regulierungen reduzieren den Vendor Lock-In in wichtigen 
Teilbereichen: Sie schaffen Wechselrechte, verbieten 
Egress-Gebühren im Kontext von Anbieterwechseln und 
untersagen wettbewerbswidrige Koppelungspraktiken. 

Dennoch beseitigen sie nicht die fundamentale 
Marktasymmetrie. US-amerikanische Hyperscaler verfügen 
über Skaleneffekte, globale Infrastruktur und 
Investitionskapazitäten, die europäische Anbieter 
strukturell nicht erreichen können, unabhängig davon, 
wie gut die Regulierung ist. Hinzu kommt, dass vollständige 
digitale Souveränität in bestimmten Bereichen weder 
realistisch noch erstrebenswert ist: Bei Halbleitern 
etwa kontrolliert TSMC über 90 Prozent der weltweiten 
Fertigung fortschrittlicher Chips, und Nvidias Dominanz 
im KI-Chip-Markt ist vergleichbar ausgeprägt. (probably microsoft in saas worth noting)

Eine 
europäische Souveränität auf diesen Ebenen zu erzwingen 
würde massive Investitionen erfordern, die in keinem Verhältnis 
zum Nutzen stehen. Vollständige Souveränität ist damit 
keine sinnvolle politische Zielsetzung. Angestrebtes 
Ziel muss stattdessen eine \textit{selektive Souveränität} 
in strategisch kritischen Bereichen sein.

Der Data Act adressiert die Portabilität von Daten,
nicht jedoch die Portabilität von Anwendungen. Ein 
Unternehmen, das tief in proprietäre Cloud-Dienste 
wie AWS Lambda, Google BigQuery oder Azure Cognitive 
Services integriert ist, kann seine Daten zwar 
standardkonform exportieren. Die darauf aufbauende 
Anwendungslogik muss jedoch vollständig neu entwickelt 
werden, da diese Dienste keine standardisierten 
Äquivalente bei anderen Anbietern haben. Zudem gilt 
das Egress-Verbot ausschließlich im Kontext eines 
vollständigen Anbieterwechsels. Laufende 
Datentransferkosten im Cloud-Betrieb bleiben weiterhin 
zulässig.

Die rechtliche Exposition gegenüber dem US CLOUD Act 
bleibt durch den Data Act strukturell ungelöst. Zwar 
enthält Artikel 32 des Data Acts Schutzklauseln gegen 
den Zugriff von Drittstaatsbehörden auf 
nicht-personenbezogene Daten. Personenbezogene Daten 
fallen weiterhin ausschließlich unter die DSGVO, sodass 
der in Kapitel \ref{ch:problem} beschriebene 
Jurisdiktionskonflikt fortbesteht. Entscheidend ist 
dabei: Der CLOUD Act greift auf Ebene des 
US-Mutterkonzerns, nicht auf Ebene der EU-Tochtergesellschaft. 
Solange AWS, Microsoft und Google US-amerikanische 
Unternehmen sind, unterliegen sie dem CLOUD Act, 
unabhängig davon, was europäisches Recht vorschreibt 
und unabhängig davon, ob ihre Daten auf europäischen 
Servern gespeichert sind.

Schließlich ist die Regulierungsstrategie der EU ein 
zweischneidiges Schwert. Einerseits zeigt sie messbare 
Wirkung: Da digitale Souveränität zu einem zentralen 
politischen Thema geworden ist, haben die Hyperscaler 
begonnen, sich aktiv um Compliance zu bemühen und 
souveräne Cloud-Angebote zu entwickeln (cite souveräne cloud angebote). 
Andererseits birgt intensive Regulierung das Risiko, 
Innovation und die Wettbewerbsfähigkeit der digitalen 
Transformation in der EU zu hemmen (cite gipfel für digitale souveränität) und erzeugt, 
wie in Kapitel (cite problem) gezeigt, geopolitischen 
Gegendruck: Die US-Regierung hat die europäische 
Digitalregulierung explizit als diskriminierend gegenüber 
US-amerikanischen Technologieunternehmen bezeichnet und 
mit Gegenmassnahmen gedroht (cite trump). 
Regulierung schützt also einerseits die digitale 
Souveränität und gefährdet sie andererseits, indem 
sie zum Hebel geopolitischer Auseinandersetzungen wird.

\section{Lösungsansatz}
\label{ch:loesungsansatz}

Aus der vorangegangenen Analyse ergibt sich, dass keine 
einzelne Maßnahme die digitale Souveränität der EU 
gegenüber Cloud-Abhängigkeiten vollständig wiederherstellen 
kann. Ein realistischer Lösungsansatz muss daher mehrere 
komplementäre Strategien kombinieren, mit dem Ziel nicht 
vollständiger Unabhängigkeit, sondern verwalteter 
Abhängigkeit mit glaubwürdigen Alternativen.

\subsection{Strukturelle Trennung als Maximalmodell: 
Das Beispiel China} (cite china cloud sovereignty)

Das radikale Modell digitaler Cloud-Souveränität hat 
China realisiert: Durch gesetzliche Vorschriften wurde 
ausländischen Cloud-Anbietern der direkte Betrieb von 
Infrastruktur auf chinesischem Boden untersagt. 
US-amerikanische Anbieter wie AWS und Microsoft mussten 
ihre Technologie an chinesische Partner lizenzieren und 
die operative Kontrolle vollständig abgeben. Amazon 
besitzt in China weder die Hardware noch hat es Zugriff 
auf die dort gespeicherten Daten. 

Das Ergebnis ist eine genuine strukturelle Souveränität: 
Kein CLOUD-Act-Beschluss kann auf Daten zugreifen, über 
die Amazon operativ keine Kontrolle mehr hat. Dieses 
Modell ist für die EU aus mehreren Gründen nicht 
direkt übertragbar: Es erfordert massive staatliche 
Investitionen in heimische Alternativen, schränkt den 
Wettbewerb ein und würde erheblichen geopolitischen 
Widerstand auslösen. Es verdeutlicht jedoch, dass 
echte strukturelle Souveränität möglich ist, wenn 
der politische Wille vorhanden ist, sie durch 
Gesetz durchzusetzen.

\subsection{Bevorzugung europäischer Cloud-Anbieter 
für kritische Infrastruktur}

Unterhalb der chinesischen Maximallösung liegt ein 
pragmatischerer Ansatz: Die EU und ihre Mitgliedstaaten 
könnten gesetzlich vorschreiben, dass bestimmte 
Kategorien kritischer Infrastruktur wie z.B öffentliche 
Verwaltung, Gesundheitswesen, Verteidigung, 
Justizsysteme ausschließlich bei europäischen 
Cloud-Anbietern betrieben werden dürfen. Europäische 
Anbieter wie OVHcloud, Hetzner oder IONOS sind für 
diese Standardanwendungsfälle vollwertige Alternativen 
und unterliegen keiner fremden Rechtsordnung. Der 
ECFR empfiehlt genau diesen Ansatz: Europäische 
Technologien müssen das Rückgrat der kritischsten 
staatlichen Funktionen bilden -- auch wenn für 
weniger sensible Bereiche weiterhin US-Hyperscaler 
genutzt werden können (cite ecfr2025eurostack). 
Dieser Ansatz entspricht dem Konzept der \textit{selektiven 
Souveränität}: Nicht alles muss souverän sein,
aber das Kritischste muss es sein.

\subsection{Hybrid Cloud und Open-Source-Technologien 
zur Reduktion von Lock-In}

Für Organisationen, die auf die Leistungsfähigkeit 
der Hyperscaler angewiesen sind, bietet ein 
\textit{Hybrid-Cloud-Ansatz} eine praktische 
Zwischenlösung: Kritische und sensible Daten werden 
auf europäischer oder privater Infrastruktur betrieben, 
während weniger sensitive Workloads auf Public-Cloud-Diensten 
laufen. Ergänzt wird dies durch den Einsatz von 
agnostischen Open-Source-Technologien wie Kubernetes, OpenStack 
oder Apache Kafka, die keine Bindung an einen 
spezifischen Anbieter erzeugen und einen späteren 
Wechsel erheblich erleichtern. Wer proprietäre 
Anbieter-Dienste durch standardisierte Open-Source-Äquivalente 
ersetzt, behält die Anwendungsportabilität, die der 
Data Act (cite data act) allein nicht gewährleisten kann.

\subsection{Clientseitige Verschlüsselung 
für kritische Daten}

Eine ergänzende technische Maßnahme, die unabhängig 
vom Anbieter wirksam ist, ist die \textit{clientseitige 
Verschlüsselung}: Daten werden verschlüsselt, bevor 
sie in die Cloud übertragen werden, und der 
Verschlüsselungsschlüssel verbleibt ausschließlich 
beim Kunden, nicht beim Cloud-Anbieter. Dies bedeutet: 
Selbst wenn ein US-Gericht AWS per CLOUD Act zur 
Herausgabe von Daten verpflichtet, erhält die Behörde 
ausschließlich unlesbare verschlüsselte Daten. 
Clientseitige Verschlüsselung löst das 
Jurisdiktionsproblem nicht auf rechtlicher Ebene.
Sie macht es jedoch auf technischer Ebene weitgehend 
irrelevant. Für besonders sensitive staatliche und 
institutionelle Daten ist dies eine der wirksamsten 
und unmittelbar umsetzbaren Schutzmaßnahmen.
